% Appendix S: the six-step derivation, from the horizon to the cell (the saturation-appendix task brief, part 2).
% Working title "Saturation: one cost from the horizon to the cell". THE AUTHOR PICKS THE FINAL TITLE.
% Steps 1-5 carry the author's own text (the source text for Steps 1-5, from the author's derivation paper, 2026-10-05) near word for
% word, with the book's equations and macros. Changes from the source, each to agree with the book as it stands: one nat (not one bit) per
% 4 l_P^2 (PAPER_ERRATA T9; eq:law_dAmin); "floor" at the black hole read as the ceiling (brief: the horizon is the ceiling, H_min is the floor);
% the CpG count is the hg19 index the book uses (Ch. ch:landauer), not the source's 19.6e6; "37 orders of magnitude" for the parameters replaced
% by the computed ratios (Ch. ch:bridge uses 37 orders for the span of length); the modified Friedmann equation written for the matter sector
% (eq:Hm), with the background unchanged; E(a) named the activation function, as in the book. The source's empirical-confirmation section
% (pre-chain results) is left out. Step 6 is written new from chain v3 (Biological_Physics/MethylPhys/chain, sop/MethylPhys_CPG_SOP_v3.md)
% and the Part VI chapters. The Hawking temperature is written T_BH, as in Ch. ch:blackholes, because T_H is the cosmic horizon in Step 5.
% Every number is checked in docs/book/verify_book.py (labels app:saturation:L<line>).
\chapter{Saturation: one cost from the horizon to the cell}\label{app:saturation}

A system reaches the limit of what it can maintain when its rate of irreversible information production saturates the capacity of its
nearest encoding surface. At stellar scale this is the black hole. At cell scale it is the cell's record leaving its maintained state.

This appendix sets out the chain that joins the two, in six steps. Four are established physics: Bekenstein, Hawking, Landauer and Jacobson. The
fifth is IAM's one new identification. The sixth carries the same cost to the cell. Each step names the chapter that develops it, and every
equation carries its status label. The chain is about encoding surfaces and what they can hold. It does not change the cosmic background: in IAM
the background expansion is the standard one, and the record term enters the matter sector only (Chapter~\ref{ch:iams_law},
Section~\ref{sec:law_caikim}). The terms used here (encoding surface, floor, full surface, Met-A, IAM-A) are in the glossary,
Appendix~\ref{app:glossary}.

\section{The floor and the ceiling}\label{app:sat:ends}
Every surface in this appendix has two ends, and they are kept apart throughout.
\begin{itemize}
\item \textbf{The floor} is $\Hmin$, the minimum entropy a healthy maintained record holds. It is measured on purified healthy cells
      (Step~6, Section~\ref{app:sat:step6}). A horizon has no floor.
\item \textbf{The ceiling} is the surface saturated at its full capacity. The black hole is this end: a horizon written to the
      Bekenstein--Hawking bound. A collapsing star approaches its ceiling. For the cell the ceiling is the full disorder of its record, every
      site at a coin flip (Chapter~\ref{ch:surface}, Section~``A capacity, and a ceiling'').
\end{itemize}

\section{Step 1 --- Bekenstein (1973): horizon entropy}\label{app:sat:step1}
Bekenstein showed that a black hole of horizon area $A$ carries entropy proportional to that area~\cite{Bekenstein1973}. The
Bekenstein--Hawking entropy is
\begin{equation}
S_{\rm BH}=\frac{\kB A}{4\lP^2},\qquad \lP=\sqrt{\frac{\hbar G}{c^3}}, \label{eq:sat_SBH}
\end{equation}
where $\lP$ is the Planck length. \derived{} (established physics; Chapter~\ref{ch:blackholes}, Section~\ref{sec:bh_thermo})

The minimum encoding area is $4\lP^2$. The horizon holds one unit of entropy, one nat, per four Planck areas, and one bit per
$4\ln2\,\lP^2=2.77\,\lP^2$ (Eq.~\eqref{eq:law_dAmin}). This is not a coincidence: the coefficient $1/4$ is derived from the ratio
\begin{equation}
\eta=\frac{1}{4\lP^2},\qquad \frac14=\frac{2\pi}{8\pi}, \label{eq:sat_eta}
\end{equation}
where $2\pi$ is the Euclidean Rindler cone angle and $8\pi$ is the Einstein equation normalization. \derived{} (Chapter~\ref{ch:bekenstein},
Eq.~\eqref{eq:bk_eta}) The encoding surface has a definite capacity. When that capacity is saturated, the stellar system has reached its
ceiling.

\section{Step 2 --- Hawking (1975): the encoding surface has a temperature}\label{app:sat:step2}
Hawking showed that black holes radiate thermally~\cite{Hawking1975}, at the Hawking temperature of a black hole of mass $M$:
\begin{equation}
T_{\rm BH}=\frac{\hbar c^3}{8\pi GM\kB}. \label{eq:sat_TBH}
\end{equation}
\derived{} (established physics; Chapter~\ref{ch:blackholes}, Eq.~\eqref{eq:bh_Tuniv} with $\kappa=c^4/4GM$)

This is the temperature at which the encoding surface operates. Landauer's principle now applies: each irreversible bit written to this surface
costs $\kB T_{\rm BH}\ln2$ in free energy. The encoding surface is not passive storage. It is an active thermodynamic boundary with a cost per
operation.

\section{Step 3 --- Landauer (1961): the cost per bit is universal}\label{app:sat:step3}
Landauer's principle states that the irreversible erasure of one bit of information in a physical system at temperature $T$ dissipates a minimum
energy~\cite{Landauer1961}
\begin{equation}
E_{\rm bit}=\kB T\ln2 . \label{eq:sat_Ebit}
\end{equation}
\derived{} The bound has been approached in experiment~\cite{Berut2012}. \observed{} (Chapter~\ref{ch:iams_law}, Eq.~\eqref{eq:law_landauer};
Chapter~\ref{ch:bridge}, Eq.~\eqref{eq:landauer})

This is scale-independent. $T$ is the local temperature of the encoding surface. $N$ is the number of bits being written. The total Landauer cost
is $N$ times $E_{\rm bit}$:
\begin{equation}
E_N=N\kB T\ln2 . \label{eq:sat_EN}
\end{equation}
\derived{} At the black-hole horizon: $T=T_{\rm BH}$, $N=\delta A/(4\ln2\,\lP^2)$. At the cell nucleus: $T=\Tb=310.15$\,K, $N$ the CpG sites of
the record (Step~6). The equation is the same. The parameters differ by many orders of magnitude (Table~\ref{tab:sat:identity}).

\section{Step 4 --- Jacobson (1995): gravity is the thermodynamics of encoding surfaces}\label{app:sat:step4}
Jacobson's 1995 result is the pivot of the entire chain~\cite{Jacobson1995}. He showed that Einstein's field equations are not a fundamental
law: they are the thermodynamic equation of state for spacetime, derivable from the Clausius relation applied to local Rindler horizons,
\begin{equation}
\delta Q=T\,\delta S\quad\Longrightarrow\quad
R_{ab}-\tfrac12Rg_{ab}+\Lambda g_{ab}=\frac{2\pi}{\hbar c\,\eta}T_{ab}\equiv\frac{8\pi G}{c^4}T_{ab},\qquad G=\frac{c^3}{4\hbar\eta}. \label{eq:sat_jacobson}
\end{equation}
\derived{} given $\delta Q=T\,\delta S$ on every local causal horizon (Assumption~1, Section~\ref{sec:assumptions}). The derivation step by step is in
Chapter~\ref{ch:iams_law}, Eq.~\eqref{eq:law_jacobson}, Chapter~\ref{ch:bekenstein}, Section~\ref{sec:bk_jacobson}, and
Appendix~\ref{app:derivations}, Section~\ref{app:der:jacobson}.

The entropy functional $S$ determines the theory of gravity. Jacobson used $S_{\rm geo}=\kB\eta A$ alone, with $\eta=1/(4\lP^2)$. The result is
general relativity.

This means gravity is not a force that exists alongside thermodynamics. Gravity is the thermodynamic response of spacetime to information being
written to encoding surfaces. The encoding surface is not a consequence of gravity. Gravity is a consequence of the encoding surface. \interp

\section{Step 5 --- IAM: adding $S_{\rm info}$ to the entropy functional}\label{app:sat:step5}
The single new physical identification in IAM: decoherence-produced classical information contributes to horizon entropy. The total entropy
functional becomes
\begin{equation}
S_{\rm total}=S_{\rm geo}+S_{\rm info}, \label{eq:sat_Stotal}
\end{equation}
where $S_{\rm info}$ is the accumulated Landauer cost of every irreversible quantum-to-classical transition in the bulk since the electroweak
transition. \conjecture{} (the one new identification; Chapter~\ref{ch:iams_law}, Eq.~\eqref{eq:law_Stotal})

Applying the Cai--Kim extension of the Jacobson mechanism~\cite{CaiKim2005} to the FRW apparent horizon with this extended entropy functional
yields the modified Friedmann equation of the matter sector:
\begin{equation}
H_m^2=\frac{8\pi G}{3}\rho_m+\frac\Lambda3+\beta_mE(a)H_0^2 , \label{eq:sat_Hm}
\end{equation}
\derived{} given Assumptions~1 and~2 (Chapter~\ref{ch:iams_law}, Section~\ref{sec:law_caikim}, Eq.~\eqref{eq:Hm}). The term enters the
matter perturbations only; the background expansion and every photon-sector observable are unchanged. That placement is the sector rule
of Chapter~\ref{ch:iams_law} (Section~\ref{sec:law_dualsector}). \conjecture

Here $E(a)=\exp(1-1/a)$ is the activation function, from the record constraint integrated over cosmic time (Eq.~\eqref{eq:Ea}).
\derived{} given the record constraint. The coupling is the virial share of the matter density,
\begin{equation}
\beta_m=\frac{\Omega_m}{2}=0.15765\quad(\Omega_m=0.3153), \label{eq:sat_beta}
\end{equation}
\prediction{} (Eq.~\eqref{eq:law_beta}) with zero free parameters; it is fixed before any chain runs, and the Planck~2018 Level~2 chain
posterior is consistent with it at $0.2\sigma$ (Chapter~\ref{ch:dual}). \measured{} $\Lambda$ is the cosmological constant, which IAM reads as
the vacuum baseline of irreversible records (Chapter~\ref{ch:lambda}). \conjecture

Black holes in IAM are mandatory local encoding surfaces: when the local decoherence rate of a collapsing stellar system saturates the
Bekenstein--Hawking bound at the local horizon, the system reaches its ceiling. It can no longer pay the cost of remaining a star.

\begin{keybox}[title=Black hole: the saturation condition (stellar scale)]
A star reaches its ceiling when
\begin{equation}
\frac{dS_{\rm local}}{dt}\ \ge\ \frac{dS_{\rm BH}}{dt}. \label{eq:sat_star}
\end{equation}
The local decoherence rate saturates the Bekenstein--Hawking encoding capacity. The system can no longer maintain its thermodynamic identity as
a star. The encoding surface absorbs what the stellar system can no longer maintain. $T=T_{\rm BH}$; $N=A_{\rm horizon}/(4\ln2\,\lP^2)$ bits.
\conjecture{} (Chapter~\ref{ch:blackholes}, Section~\ref{sec:bh_formation}, Eq.~\eqref{eq:bh_saturation}; at the Schwarzschild radius the
bound is met exactly, Eq.~\eqref{eq:bh_hoop})
\end{keybox}

\section{Step 6 --- The cell: the methylation record}\label{app:sat:step6}
At cell scale the encoding surface is the methylation record: the pattern of methylated and unmethylated CpG sites the cell writes on its
chromatin and copies at every division (Chapter~\ref{ch:surface}). Each site is a binary record with a capacity of one bit
(Eq.~\eqref{eq:H}). The surface operates at body temperature, $\Tb=310.15$\,K, and the cost per bit is
\begin{equation}
E_{\rm bit}=\kB\Tb\ln2=2.97\times10^{-21}\ {\rm J}. \label{eq:sat_ebit}
\end{equation}
\calc{} (Chapter~\ref{ch:landauer}, Eq.~\eqref{eq:ebit})

$N$ is the number of CpG sites in the record. Counting every CpG of the hg19 CpG index the book uses gives $N=28{,}217{,}448$
(Chapter~\ref{ch:landauer}, Eq.~\eqref{eq:efloor}). \calc{} The cost of the whole record is Eq.~\eqref{eq:sat_EN} at $T=\Tb$: the same
equation as at the horizon.

\paragraph{The floor.} The floor is $\Hmin$, the minimum entropy a healthy maintained record holds. It is measured on purified healthy cells,
through chain v3, with nothing fitted (Chapter~\ref{ch:chain}). On single molecules its form is the two-state Boltzmann floor of the copy
error, $\Hmin=H(\varepsilon_0)$, with its height from the holding energy measured on healthy molecules (Chapter~\ref{ch:iama},
Eq.~\eqref{eq:eps0}). \derived{} for the form, \measured{} for the height.

\paragraph{The reading.} Met-A on arrays (Chapter~\ref{ch:meta}) and IAM-A on single-molecule sequencing (Chapter~\ref{ch:iama}) read how far a
cell's record sits above that floor. Each reads the cell against purified healthy cells of the same type, on the same instrument, so that a
healthy cell reads 1 on the one gauge of Chapter~\ref{ch:gauge}. \calibrated{} (the Met-A reference) \measured{} (the IAM-A position)

\paragraph{The ceiling.} The ceiling is the full disorder of the record: every site at $\beta=\tfrac12$, one bit each, the full surface of
Chapter~\ref{ch:floorbreach}. \derived{} The cell approaches it through a loss of maintained fidelity, not through a shortage of energy
(Chapter~\ref{ch:landauer}).

\paragraph{The identification.} The cell's approach to its ceiling is identified with the black hole's approach to its ceiling: the same
cost per bit, the same notion of a surface written to capacity, at a different temperature and a different number of bits. \conjecture

\section{The structural identity}\label{app:sat:identity}
Table~\ref{tab:sat:identity} sets the two ends of the chain side by side.
The same equation, $N\kB T\ln2$, stands at both ends. $T$ and $N$ differ by many orders of magnitude: the temperatures by a factor of
$5.0\times10^{9}$ and the bit counts by $5.4\times10^{69}$. \calc

\begin{table}[htbp]\centering\small
\caption{Black hole (one solar mass) against cell (hg19 CpG index). \calc{} for the numbers; the identification of the two ceilings is
\conjecture.}\label{tab:sat:identity}
\begin{tabularx}{\textwidth}{@{}>{\raggedright\arraybackslash}p{0.2\textwidth}>{\raggedright\arraybackslash}X>{\raggedright\arraybackslash}X@{}}\toprule
 & black hole & cell \\\midrule
encoding surface & the horizon & the methylation record (CpG sites on chromatin) \\\rowrule
temperature & $T_{\rm BH}=\hbar c^3/8\pi GM\kB\approx6.17\times10^{-8}$\,K & $\Tb=310.15$\,K \\\rowrule
bit count & $A/4\lP^2\approx1.05\times10^{77}$ nats, $1.51\times10^{77}$ bits & $N$ sites, $2.82\times10^{7}$ \\\rowrule
cost per bit & $\kB T_{\rm BH}\ln2\approx5.90\times10^{-31}$\,J & $\kB\Tb\ln2=2.97\times10^{-21}$\,J \\\rowrule
floor & none & $\Hmin$, measured on purified healthy cells \\\rowrule
ceiling & horizon saturation, $S=S_{\rm BH}$ & full disorder of the record, every site at $\beta=\tfrac12$ \\\rowrule
what is read & the horizon & Met-A (arrays), IAM-A (single molecules) \\\bottomrule
\end{tabularx}
\end{table}
